\subsubsection*{命题 1.2.}
对 $p$ 的任意幂 $q$, 若
\begin{equation*}
\sum_{i=1}^{r}d_i\geq n+r+1,
\end{equation*}
则在 $\mathbb P^{n+r}_{\mathbb F_q}$ 中存在多重次数为 $(d_1,\dots,d_r)$
且没有 $\mathbb F_q$ 有理点的完全交.

\medskip
\noindent\emph{证明.}
先考虑 $r=1$ 的情形. 对每个 $d\geq n+2$, 需找到次数为 $d$ 的齐次式
\begin{equation}\tag{1.2.1}
H(X_1,\dots,X_{n+2})\in
\mathbb F_q[X_1,\dots,X_{n+2}],
\end{equation}
使得
\begin{equation}\tag{1.2.2}
\begin{aligned}
&(\alpha_1,\dots,\alpha_{n+2})\in(\mathbb F_q)^{n+2},\qquad
H(\alpha_1,\dots,\alpha_{n+2})=0\quad\Longrightarrow\\
&\hspace{7em}\alpha_1=\cdots=\alpha_{n+2}=0.
\end{aligned}
\end{equation}
置 $\mathbb E=\mathbb F_{q^d}$, 并选取 $\mathbb E$ 中在 $\mathbb F_q$ 上线性无关的
元素 $\beta_1,\dots,\beta_{n+2}$. 取
\begin{equation}\tag{1.2.3}
H(X_1,\dots,X_{n+2})=
N_{\mathbb E/\mathbb F_q}
\left(\sum_{i=1}^{n+2}\beta_iX_i\right).
\end{equation}

在一般情形, 由于
\begin{equation}\tag{1.2.4}
\sum_{i=1}^{r}d_i\geq n+r+1,
\end{equation}
可以置 $a_0=0$, 并选取 $r$ 个整数
\begin{equation}\tag{1.2.5}
1\leq a_1<a_2<\cdots<a_r=n+r+1
\end{equation}
使得
\begin{equation}\tag{1.2.6}
d_i\geq a_i-a_{i-1}\qquad(1\leq i\leq r).
\end{equation}
对每个 $i=1,\dots,r$, 选取次数为 $d_i$ 的齐次式
\begin{equation*}
H_i(X_{a_{i-1}+1},\dots,X_{a_i})
\in\mathbb F_q[X_{a_{i-1}+1},\dots,X_{a_i}],
\end{equation*}
其变量为所列的 $a_i-a_{i-1}$ 个变量, 且具有性质 (1.2.2).
所求完全交由下列齐次式定义
\begin{equation*}
H_1(X_1,\dots,X_{a_1}),\quad\dots,\quad
H_r(X_{a_{r-1}+1},\dots,X_{a_r}).
\end{equation*}

\medskip
\noindent\emph{底本注.}
照录法文层保留了印本 (1.2.1) 中错置的括号、(1.2.6) 中字面的界 $d_1\geq a_o$,
以及末尾 $H_r$ 块中漏去的 $+1$. 对原始像素的直接检查确认这些确为印本特征;
此构造要求采用上列相容形式.
